Os ensaios foram planejados para avaliar a camada WCAN isoladamente, sem tornar
o barramento CAN físico o fator dominante da análise. A integração com TWAI foi
validada previamente em bancada, mas a campanha final concentrou-se no enlace
ESP-NOW, nas filas internas, na filtragem de identificadores CAN e nos mecanismos
de confirmação adotados pelas variantes \emph{broadcast} e \emph{multicast}. O
objetivo metodológico foi produzir uma campanha reprodutível, com variação
controlada de topologia, frequência, transporte e política de agrupamento
temporal.

A bancada experimental foi organizada em torno de um computador Ubuntu usado para
orquestrar os ensaios, gravar os \emph{firmwares}, sincronizar as placas e coletar
os \emph{logs} seriais. A campanha final utilizou cinco placas conectadas por
USB: três ESP32, identificadas como A, B e C, e duas ESP32-C3, identificadas como
D e E. As placas ESP32 podiam atuar como sensores ou receptores, enquanto as
ESP32-C3 foram usadas apenas como sensores. Essa restrição foi adotada porque o
receptor concentra recepção, filtragem, deduplicação e emissão de métricas, papel
mais adequado às placas ESP32 usadas na bancada, que possuem dois núcleos. Cada
ensaio teve duração nominal de 30 s, com 5 s de intervalo entre execuções e
comunicação serial a 921600 bit/s.

A matriz final considerou topologias com pelo menos um sensor e um receptor,
respeitando o limite de cinco placas disponíveis e a restrição de que apenas as
três ESP32 podiam ser receptoras. Assim, o plano incluiu as topologias 1S-1R,
1S-2R, 1S-3R, 2S-1R, 2S-2R, 2S-3R, 3S-1R, 3S-2R e 4S-1R. Para as suítes com
frequência fixa, foram avaliados 10, 100, 200, 750 e 1000 Hz. Na suíte de
frequências mistas, cada sensor recebeu uma frequência distinta escolhida entre
100, 200, 750 e 1000 Hz. As cinco suítes executadas estão resumidas na
Tabela~\ref{tab:suites-executadas}; em conjunto, elas consolidaram a campanha
final em 684 execuções planejadas.

\begin{table}[htbp]
\centering
\scriptsize
\setlength{\tabcolsep}{3pt}
\renewcommand{\arraystretch}{1.16}
\begin{tabular}{p{0.17\linewidth}|p{0.24\linewidth}|p{0.27\linewidth}|p{0.22\linewidth}}
\hline
Suíte & Configuração & Topologias e frequências & Objetivo metodológico \\
\hline
\texttt{baseline} & 1 identificador CAN por sensor; \texttt{linger\_ms = 100}. & Matriz completa; 10, 100, 200, 750 e 1000 Hz. & Referência com agrupamento temporal. \\
\hline
\texttt{multiple} & 3 identificadores CAN por sensor; \texttt{linger\_ms = 100}. & Matriz completa; 10, 100, 200, 750 e 1000 Hz. & Concorrência entre fluxos do mesmo sensor. \\
\hline
\texttt{real\_time} & 1 identificador CAN por sensor; \texttt{linger\_ms = 0}. & Matriz completa; 10, 100, 200, 750 e 1000 Hz. & Envio amostra a amostra. \\
\hline
\texttt{active\_filter} & 1 identificador CAN por sensor; \texttt{linger\_ms = 100}; filtros por receptor. & Topologias com pelo menos dois sensores; 10, 100, 200, 750 e 1000 Hz. & Entrega apenas a receptores interessados. \\
\hline
\texttt{mixed\_frequency} & 1 identificador CAN por sensor; \texttt{linger\_ms = 100}; frequências distintas por sensor. & Topologias com pelo menos dois sensores; 100, 200, 750 e 1000 Hz mistos. & Convivência de frequências distintas. \\
\hline
\end{tabular}
\caption{Suítes executadas na campanha final e objetivo metodológico de cada uma.}
\label{tab:suites-executadas}
\end{table}

A taxa de entrega foi calculada a partir dos contadores gerados pelos sensores e
dos intervalos recebidos pelos receptores. Para cada sensor e identificador CAN,
o analisador reconstruiu o conjunto de contadores transmitidos e o comparou com
a união dos conjuntos recebidos pelos receptores esperados. O critério principal
de entrega adotado na campanha foi, portanto, o recebimento por pelo menos um
receptor esperado. A análise par a par entre um sensor e cada receptor individual
foi preservada como informação diagnóstica, mas não como critério principal de
aprovação. Na suíte \texttt{active\_filter}, apenas os receptores cuja
\emph{allowlist} contém determinado identificador CAN entram no conjunto de
receptores esperados para aquele fluxo.

Além da entrega, foram coletadas falhas observadas no caminho de transmissão,
comportamento de agrupamento temporal, frequência média observada, estatísticas
de lote, tempo estimado de \emph{airtime}, falhas de envio do sensor e métricas
de memória \emph{heap}. O analisador também verificou condições de validade do
ensaio, como encerramento normal da execução, abortos, mensagens de \emph{panic}
ou \emph{reset}, cumprimento dos filtros da suíte \texttt{active\_filter} e
compatibilidade entre a frequência observada e a frequência configurada. Para as
frequências fixas e mistas, a tolerância definida para essa comparação foi de
5\%. Os arquivos produzidos pela campanha, os detalhes de automação, o protocolo
UART e a estrutura dos \emph{logs} são registrados no
Apêndice~\ref{appendix:full-run}.
